\begin{preface}{Abstract}
  The advent of machine learning and edge computing has revolutionized motion tracking applications, enabling real-time, energy-efficient, and accurate systems for diverse use cases in healthcare, sports analytics, and human-computer interaction. However, developing AI models for motion tracking remains complex, requiring specialized expertise in data preprocessing, feature extraction, model training, and edge deployment. This paper presents a comprehensive framework that simplifies the end-to-end process of generating AI models tailored for human motion tracking on resource-constrained edge devices.

  The framework introduces a novel \textit{interactive preprocessing approach} featuring draggable time window selection, enabling users to efficiently annotate and segment motion data without extensive programming knowledge. The system supports multiple machine learning algorithms including Random Forest, Support Vector Machines, and Neural Networks, with automated code generation for various edge platforms including Arduino, ARM Cortex-M, and Seeed XIAO microcontrollers. Our optimization-aware code generation considers multiple deployment strategies—accuracy, speed, power consumption, and balanced approaches—adapting the generated code to specific application requirements.

  We validate the framework using a Human Activity Recognition (HAR) dataset collected at 100Hz from a 6-axis IMU sensor (LSM6DS3) on the Seeed XIAO nRF52840 Sense platform. In this proof-of-concept study using single-subject data across five activity classes, experimental results demonstrate that deployed models achieve 94.5\% accuracy for Random Forest and 96.2\% for Neural Networks, with inference times of 15ms and 12ms respectively, while consuming less than 10mW during operation. Comparative analysis with commercial platforms (Edge Impulse, SensiML) reveals that our open-source framework offers superior preprocessing flexibility and code generation quality at zero cost.

  The framework addresses critical gaps in current solutions by providing: (1) an intuitive GUI-based workflow accessible to non-experts, (2) complete control over the preprocessing pipeline with visual feedback, (3) hardware-agnostic deployment with platform-specific optimizations, and (4) an extensible architecture supporting future integration of additional algorithms and platforms. This research contributes to democratizing edge AI development for motion tracking applications, with potential impact in healthcare rehabilitation, athletic performance monitoring, and assistive technologies.
\end{preface}
